\chapter{Groups as Symmetries}

\begin{center}\itshape
A cube can be put back in its box in $24$ ways, and shuffling four cards can be done in $24$ ways.\\
These are the same $24$, in a sense this chapter makes precise.
\end{center}

We have already met groups: the roots of unity in Volume~1 and the invertible residues modulo $m$ in Volume~2. Here we stop treating their rules as special facts about numbers. We ask which rules work for rotations, permutations and matrices.

Put a square on the table and mark one corner. Turn it through a quarter-turn. The square still fills exactly the same space, but the mark has moved. Turn again, and the two turns together are equivalent to one half-turn. Four quarter-turns bring the mark home.

Now allow reflections as well. The square has eight symmetries: four rotations and four reflections. They can be combined, but the order can matter. Reflect first and rotate second; you may get a different result from rotating first and then reflecting. The set of eight moves, with the rule for combining them, is the group $D_4$.

Integers under addition, non-zero real numbers under multiplication, and
permutations under composition have the same basic rules: combine, leave
unchanged, undo. We will state those rules precisely and use them to prove
facts about all three examples.

A cube offers a first counting question. Choose the face that goes on top
($6$ choices), then which of its four edges faces the front ($4$ choices).
Those choices fix the orientation, giving $24$ rotations. Later we will see
why these rotations can be identified with the permutations of four long
diagonals. First we need to decide what counts as a group, what its smaller
groups look like, and why the size of a subgroup divides the size of a finite group.

The distinction between a rotation and its description matters. Two strings of moves can produce the same result; the group remembers the result, not the instructions used to make it. That observation will become the kernel of a homomorphism in Chapter~2.

\section{Groups}

\begin{definition}\label{def:group}
A \emph{group} is a set $G$ with a map $G\times G\to G$, $(a,b)\mapsto ab$, such that
\begin{enumerate}[label=(\roman*)]
\item (associativity) $(ab)c=a(bc)$ for all $a,b,c\in G$;
\item (identity) there is $e\in G$ with $ea=ae=a$ for all $a\in G$;
\item (inverses) every $a\in G$ has some $b\in G$ with $ab=ba=e$.
\end{enumerate}
The group is \emph{abelian} if $ab=ba$ for all $a,b$. Its \emph{order} $|G|$ is its number of elements.
\end{definition}

For abelian groups the operation is often written $a+b$, the identity $0$ and the inverse $-a$.

\begin{proposition}[first rules]\label{prop:rules}
In a group $G$:
\begin{enumerate}[label=(\alph*)]
\item the identity is unique, and each $a$ has a unique inverse, written $a^{-1}$;
\item (cancellation) $ab=ac$ implies $b=c$, and $ba=ca$ implies $b=c$;
\item $(ab)^{-1}=b^{-1}a^{-1}$ and $(a^{-1})^{-1}=a$;
\item for $a,b\in G$ the equation $ax=b$ has exactly one solution, $x=a^{-1}b$. So in the multiplication table of a finite group, every row (and every column) contains each element exactly once.
\end{enumerate}
\end{proposition}
\begin{proof}
We will use the group axioms one at a time. Associativity matters:
it lets us change the parentheses, but it never lets us exchange the
order of two factors.

\emph{(a) Uniqueness.} If $e$ and $e'$ are both identities, then
$e=ee'=e'$. If $b$ and $c$ are both inverses of $a$, then
\[
b=be=b(ac)=(ba)c=ec=c.
\]
So neither the identity nor an inverse involves a choice.

\emph{(b) Cancellation.} From $ab=ac$, multiply on the left by
$a^{-1}$:
\[
a^{-1}(ab)=a^{-1}(ac)\quad\Longrightarrow\quad b=c.
\]
For $ba=ca$, multiply on the right by $a^{-1}$ instead. We can
cancel precisely because an inverse exists.

\emph{(c) Reversing products.} Check both multiplication orders:
\[
(ab)(b^{-1}a^{-1})=a(bb^{-1})a^{-1}=e,
\]
and
\[
(b^{-1}a^{-1})(ab)=b^{-1}(a^{-1}a)b=e.
\]
By uniqueness, $(ab)^{-1}=b^{-1}a^{-1}$. The element $a$ is an
inverse of $a^{-1}$, so $(a^{-1})^{-1}=a$ as well. Notice that
the order of the factors has reversed.

\emph{(d) Equations and tables.} The element $x=a^{-1}b$
satisfies $ax=a(a^{-1}b)=b$. Conversely, if $ax=b$, left
multiplication by $a^{-1}$ forces $x=a^{-1}b$. Thus each
equation $ax=b$ has exactly one solution. The same reasoning
on the right shows that $xa=b$ has exactly one solution
$x=ba^{-1}$. In a finite group's multiplication table, each row
and each column therefore lists every group element once.
\end{proof}

Powers are defined by $a^0=e$, $a^{n+1}=a^na$ for $n\ge0$, and
$a^{-n}=(a^{-1})^n$ for $n\ge1$. Associativity joins strings of copies
of $a$, giving $a^ma^n=a^{m+n}$ for nonnegative $m,n$. For opposite
signs, cancel adjacent pairs $aa^{-1}=a^{-1}a=e$: if $m\ge n\ge0$,
$a^m(a^{-1})^n=a^{m-n}$; if $n>m$, the remaining factors give
$(a^{-1})^{n-m}$. Reversing the roles gives the other mixed case, and
two negative powers obey the positive rule applied to $a^{-1}$.
Thus $a^ma^n=a^{m+n}$ for all integers $m,n$.
Repeating a block gives $(a^m)^n=a^{mn}$ for $n\ge0$;
for negative $n$, use $(a^m)^{-1}=a^{-m}$, which follows from the
first law. These rules concern powers of the same element; no commutation
of two different group elements is assumed.

Here a \emph{field} $K$ has an abelian addition group with identity $0$;
its multiplication is associative and commutative with identity $1$,
distributes over addition, satisfies $0\ne1$, and gives every nonzero
element a multiplicative inverse. The prime residue fields of Volume~2,
Chapter~3 are examples. Distributivity gives $c0=c(0+0)=c0+c0$;
additive cancellation therefore gives $c0=0$. If $ab=0$ and $a\ne0$,
then $b=a^{-1}(ab)=0$. Thus nonzero elements are closed under products
and form an abelian multiplication group. The coordinate arguments from Volume~1's
linear-algebra chapter apply over any such $K$: the exchange proof divides
only by a nonzero coefficient, and the determinant proof uses only finite
sums, products and cancellation. In particular, $n$ independent columns
in $K^n$ are a basis, $\det(AB)=\det A\det B$, and a square matrix is
invertible exactly when its determinant is nonzero. These facts apply over
$\mathbb F_p=\Z/p\Z$ as well as over $\R$ and $\C$; no order or notion
of length on $K$ is needed.

\begin{example}[the examples of this volume]\label{ex:examples}
\begin{enumerate}[label=(\arabic*)]
\item $(\Z,+)$, $(\Q,+)$, $(\R,+)$, $(\C,+)$, and $(\Q^*,\times)$, $(\R^*,\times)$, $(\C^*,\times)$: abelian.
\item For $n\ge1$, $\Z/n\Z$ under addition, and $(\Z/n\Z)^*$, the invertible classes, under multiplication (Volume~2, Chapter~3). The group of $n$-th roots of unity $\mathbb U_n=\{z\in\C:z^n=1\}$ (Volume~1, Chapter~5).
\item The \emph{symmetric group} $\mathfrak S(X)$ of all bijections of a set $X$, under composition: $(\sigma\tau)(x)=\sigma(\tau(x))$. For $X=\{1,\dots,n\}$ it is written $\mathfrak S_n$, and $|\mathfrak S_n|=n!$, by choosing the images successively in $n,n-1,\ldots,1$ ways. For distinct $a,b\in X$, the notation $(a\,b)$ means the permutation that exchanges $a$ and $b$ and fixes everything else. The group is not abelian for $n\ge3$: $(1\,2)(2\,3)$ sends $3\mapsto1$, while $(2\,3)(1\,2)$ sends $3\mapsto2$.
\item The \emph{general linear group} $GL_n(K)$ of invertible $n\times n$ matrices over a field $K$, and its subgroup $SL_n(K)$ of matrices of determinant $1$ (Volume~1, Chapter~7).
\item The \emph{dihedral group} $D_n$ ($n\ge3$): the isometries of the plane that map the set of vertices of a regular $n$-gon to itself.
\end{enumerate}
\end{example}

In (3) the product $\sigma\tau$ means ``first $\tau$, then $\sigma$''.
Associativity follows by evaluating at each $x$: both
$((\sigma\tau)\upsilon)(x)$ and $(\sigma(\tau\upsilon))(x)$ equal
$\sigma(\tau(\upsilon(x)))$. The identity map leaves every $x$ fixed,
and the inverse bijection undoes the move.

\begin{proposition}[the dihedral group]\label{prop:dihedral}
Let $n\ge3$ be an integer and $P$ a regular $n$-gon centred at $0$, with vertices $v_k=e^{2ik\pi/n}$ in $\C=\R^2$. Let $r$ be the rotation $z\mapsto e^{2i\pi/n}z$ and $s$ the reflection $z\mapsto\bar z$. Then
\[
D_n=\{e,r,r^2,\dots,r^{n-1},\ s,sr,sr^2,\dots,sr^{n-1}\},
\]
these $2n$ elements are distinct, $r^n=s^2=e$, and $srs=r^{-1}$.
\end{proposition}
\begin{proof}
\emph{A symmetry fixes the centre.}
A Euclidean isometry of the plane has the form
$f(x)=Ax+b$ with $A$ orthogonal. To see this, set
$g(x)=f(x)-f(0)$ and use the identity
\[
2\langle x,y\rangle=|x|^2+|y|^2-|x-y|^2
\]
to see that $g$ preserves scalar products. The vectors
$u_1=g(e_1),u_2=g(e_2)$ are an orthonormal basis.
For $x=x_1e_1+x_2e_2$ we have
$\langle g(x),u_i\rangle=\langle x,e_i\rangle=x_i$;
therefore $g(x)=x_1u_1+x_2u_2$. This proves linearity and
orthogonality, with $b=f(0)$. The formula $f(x)=Ax+b$
also shows that $f$ preserves averages of finitely many points.

Every polygon symmetry permutes its $n$ vertices.
Their average is zero, because
$\sum_{k=0}^{n-1}e^{2\pi ik/n}=0$ (Volume~1).
So $f(0)=0$ and $f$ is linear orthogonal.

\emph{Two adjacent vertices determine everything.}
Assume $n\ge3$, as for a polygon. The vectors
$v_0=1$ and $v_1=e^{2\pi i/n}$ are linearly
independent over $\mathbb R$. A linear symmetry
is therefore completely determined by their two images.

There are $n$ possible images $v_j$ of $v_0$.
Since an isometry preserves distance, the image of
$v_1$ must be a vertex at distance $|v_1-v_0|$
from $v_j$. Exactly two vertices have that distance:
$v_{j+1}$ and $v_{j-1}$, with indices modulo $n$.
Indeed the squared distance is
$2-2\operatorname{Re}(v_k\overline{v_j})$; fixing that real part
leaves just these two points on the unit circle, distinct because $n\ge3$.
Consequently there are at most $2n$ symmetries.

For each $j$, the rotation $r^j$ maps $(v_0,v_1)$
to $(v_j,v_{j+1})$. The map $r^js$ maps it to
$(v_j,v_{j-1})$. These are the $2n$ distinct
possibilities just counted. We may equally write
the reflections as $sr^k$: both $r^js$ and $sr^{-j}$ send $z$ to
$e^{2\pi ij/n}\bar z$, so they are the same map.
Thus the displayed list is complete.

\emph{The defining relations.}
Applying $r$ $n$ times turns through $2\pi$,
so $r^n=e$. Reflecting twice gives $s^2=e$.
Finally, for any complex $z$,
\[
srs(z)
=\overline{e^{2\pi i/n}\bar z}
=e^{-2\pi i/n}z
=r^{-1}(z).
\]
Hence $srs=r^{-1}$.

For the square, this gives four rotations and four reflections,
the eight elements in the next figure.
\end{proof}

From $srs=r^{-1}$, induction and the inverse power rules give
$sr^k=r^{-k}s$ for every integer $k$. A word has the form $r^as^\varepsilon$,
with $\varepsilon\in\{0,1\}$: appending $r$ changes $a$ to
$a+(-1)^\varepsilon$, while appending $s$ changes $\varepsilon$ to
$1-\varepsilon$. This induction on word length also gives a finite rule
for computing products. Figure~\ref{fig:symmetries} lists them for $D_4$.

\begin{think}
The face count gave $24$ cube rotations. List the possible rotation axes
and the nonzero angles that work about each. How could you check that
your list accounts for all $24$?
\end{think}

\begin{proposition}[rotations of the cube]\label{prop:cube}
The rotations of $\R^3$ that map the cube $[-1,1]^3$ to itself form a group of order $24$. Apart from the identity, they are: rotations by $\pm90^\circ$ and $180^\circ$ about the $3$ axes through the centres of opposite faces ($9$ rotations), by $\pm120^\circ$ about the $4$ long diagonals ($8$), and by $180^\circ$ about the $6$ axes through the midpoints of opposite edges ($6$).
\end{proposition}
\begin{proof}
We turn the geometric question into a finite list of matrices.

\emph{The coordinate axes must go to coordinate axes.}
Place the centre of the cube at the origin, and let $f$ be a rotation of the cube. The eight vertices $(\pm1,\pm1,\pm1)$ are exactly the points of the cube at distance $\sqrt3$ from the centre, so $f$ permutes them and sends the centre to the centre. Thus $f(0)=0$.

A map that preserves distances and fixes the origin preserves inner products. The identity already used for the polygon,
\[
2\langle x,y\rangle=|x|^2+|y|^2-|x-y|^2,
\]
has the same right-hand side for $f(x)$ and $f(y)$. The images $u_i=f(e_i)$ are therefore orthonormal. Expand $f(x)$ in that basis. For $x=x_1e_1+x_2e_2+x_3e_3$,
\[
\langle f(x),u_i\rangle=\langle f(x),f(e_i)\rangle=\langle x,e_i\rangle=x_i,
\]
so $f(x)=x_1u_1+x_2u_2+x_3u_3$. The map is linear. The matrix $A$ with these columns satisfies $A^{\mathsf T}A=I$, because the columns are orthonormal: $A$ is orthogonal. The determinant rules give $\det(A^{\mathsf T}A)=(\det A)^2=\det I=1$, so $\det A=\pm1$.

The quarter-turn about the third axis, $(x,y,z)\mapsto(-y,x,z)$, sends $e_1$ to $e_2$, $e_2$ to $-e_1$ and $e_3$ to $e_3$. The matrix with those columns has determinant $1$. Changing the sign of the first coordinate, $(x,y,z)\mapsto(-x,y,z)$, still preserves the cube and has determinant $-1$; it is a reflection, not a turn. A rotation is the distance-preserving case $\det A=1$.

Notice that
\[
e_1=\tfrac12\bigl((1,1,1)+(1,-1,-1)\bigr).
\]
Thus $Ae_1$ is the midpoint of two cube vertices,
so each of its coordinates is $-1$, $0$, or $1$.
It is also a unit vector, since $A$ is orthogonal.
The only unit vectors with those coordinates are
$\pm e_1,\pm e_2,\pm e_3$. The same argument,
with appropriate pairs of vertices, applies to
$Ae_2$ and $Ae_3$.

The three vectors $Ae_1,Ae_2,Ae_3$ form an
orthonormal basis, so they must point along
\emph{different} coordinate axes. In other words,
$A$ has precisely one non-zero entry in each
row and each column, and that entry is $\pm1$:
it is a signed permutation matrix.

\emph{Count the allowed matrices.}
There are $3!$ ways to permute the axes and
$2^3$ ways to choose the signs, giving
$3!\cdot2^3=48$ signed permutation matrices.
They all preserve the cube. Changing the sign
of the first column pairs each matrix of
determinant $1$ with a matrix of determinant
$-1$. Exactly half have determinant $1$, so there are at most
\[
\boxed{24}
\]
cube rotations. We next exhibit that many, which also shows that every determinant-$1$ signed permutation matrix in this list is a rotation.

\emph{Identify the geometric types.}
The axes through opposite face centres give
three half-turns and six quarter-turns, for $9$.
Each of the four body diagonals gives two
non-identity rotations through $120^\circ$,
for $8$. Each of the six lines through
opposite edge midpoints gives one half-turn,
for $6$. These $9+8+6=23$ non-identity
rotations are visibly distinct, and all preserve
the cube. Adding the identity gives $24$,
the full number already proved.
\end{proof}

\begin{figure}[ht]
\centering
\fig{ch01-symmetries}
\caption{Left: the multiplication table of $D_4$, the $8$ symmetries of a square ($r$ is the quarter turn, $s$ a reflection). Each row and each column contains every element once (Proposition~\ref{prop:rules}(d)). The four blocks show that rotation times rotation and reflection times reflection are rotations, while the mixed products are reflections. Right: the three kinds of rotation axes of the cube, with the number of non-trivial rotations of each kind.}
\label{fig:symmetries}
\end{figure}

The cube and the square have many symmetries, but we often want to use only some of them. The rotations of a square, for example, are closed under composition even if we forget the reflections. We call such a smaller group a subgroup.

\section{Subgroups and the order of an element}

\begin{definition}
A \emph{subgroup} of $G$ is a subset $H$ that contains $e$ and is stable under products and inverses. It is then a group for the same operation, and we write $H\le G$.
\end{definition}

\begin{proposition}\label{prop:subgroup}
\begin{enumerate}[label=(\alph*)]
\item A non-empty $H\subseteq G$ is a subgroup if and only if $ab^{-1}\in H$ for all $a,b\in H$.
\item An intersection of subgroups is a subgroup. For $S\subseteq G$, the intersection $\langle S\rangle$ of all subgroups containing $S$ is the set of finite products $s_1^{\pm1}\cdots s_k^{\pm1}$ with $s_i\in S$ (the empty product being $e$). It is the \emph{subgroup generated by $S$}.
\end{enumerate}
\end{proposition}
\begin{proof}
\emph{(a) The one-step subgroup test.}
If $H$ is a subgroup, it contains $ab^{-1}$ whenever $a,b\in H$.
Conversely, suppose $H$ is non-empty and has that property.
Choose $a\in H$; then $aa^{-1}=e\in H$. With $e\in H$, the
test applied to $e$ and any $b\in H$ shows $b^{-1}\in H$.
Finally, because $b^{-1}\in H$, the test applied to $a$
and $b^{-1}$ gives $a(b^{-1})^{-1}=ab\in H$.
We have recovered the identity, inverses, and multiplication.
So $H$ is a subgroup.

\emph{(b) Intersections and generators.}
An arbitrary intersection of subgroups is a subgroup: every
subgroup contains $e$, and if $a,b$ belong to every subgroup
in the family, then so does $ab^{-1}$. There is at least one
subgroup containing $S$, namely $G$, so the intersection
defining $\langle S\rangle$ makes sense.

Let $W$ be the set of all finite words
\[
s_1^{\varepsilon_1}s_2^{\varepsilon_2}\cdots
s_k^{\varepsilon_k},
\qquad s_i\in S,\quad\varepsilon_i\in\{-1,1\},
\]
including the empty word, whose value is $e$.
Concatenating two words gives another word.
Inverting a word reverses its letters and changes all
exponents' signs, by Proposition~\ref{prop:rules}(c).
Thus $W$ is a subgroup. It contains $S$, because a one-letter
word is allowed.

On the other hand, every subgroup containing $S$ must also
contain every inverse and every finite product of elements of
$S$, so it contains $W$. Hence $W$ is both a subgroup
containing $S$ and a subset of every such subgroup.
This proves $W=\langle S\rangle$.

For example, in the additive group $\mathbb Z$, the subgroup
generated by $4$ and $6$ is the set of all $4a+6b$ with
$a,b\in\mathbb Z$, namely $2\mathbb Z$.
\end{proof}

For instance $D_n=\langle r,s\rangle$, the rotations form the subgroup $\langle r\rangle$ of order $n$, and $SL_n(K)\le GL_n(K)$.

\begin{proposition}[order of an element]\label{prop:order}
Let $g\in G$. Either all the powers $g^k$, $k\in\Z$, are distinct, or there is a least $n\ge1$ with $g^n=e$, called the \emph{order} $o(g)$ of $g$. In the second case:
\begin{enumerate}[label=(\alph*)]
\item $g^k=e$ if and only if $n\mid k$, and $g^k=g^l$ if and only if $k\equiv l\pmod n$;
\item $\langle g\rangle=\{e,g,\dots,g^{n-1}\}$ has exactly $n$ elements;
\item $o(g^k)=\dfrac n{\gcd(n,k)}$.
\end{enumerate}
\end{proposition}
\begin{proof}
If all powers of $g$ are distinct, there is no positive $n$
with $g^n=e$: otherwise $g^n=g^0$. Conversely, if
$g^u=g^v$ for integers $u>v$, then multiplying by $g^{-v}$
gives $g^{u-v}=e$. The positive exponents for which this
happens therefore have a least member $n$.

\emph{(a) Find exactly when a power is the identity.}
Take any integer $k$, including a negative one. Euclidean
division gives $k=qn+r$ with $0\le r<n$. Hence
\[
g^k=(g^n)^qg^r=g^r.
\]
If $g^k=e$ and $r>0$, this contradicts the minimality of $n$.
Thus $r=0$, or $n\mid k$. The converse is immediate since
$g^n=e$. Finally,
\[
g^k=g^\ell\iff g^{k-\ell}=e\iff n\mid(k-\ell),
\]
which is the asserted congruence condition.

\emph{(b) Count the generated subgroup.}
Every power is one of $g^0,g^1,\ldots,g^{n-1}$ by (a).
These $n$ powers are distinct: equality of two exponents
in that range would give a non-zero multiple of $n$ whose
absolute value is less than $n$. Thus the list is exactly
$\langle g\rangle$ and has $n$ elements.

\emph{(c) Find the order of a power.}
Write $d=\gcd(n,k)$, $n=dn_0$ and $k=dk_0$,
so $\gcd(n_0,k_0)=1$. For an integer $m>0$,
\[
(g^k)^m=e
\iff n\mid km
\iff n_0\mid k_0m
\iff n_0\mid m.
\]
The last equivalence is Gauss's lemma from Volume~2:
a number coprime to $n_0$ can be cancelled in a divisibility
statement. The least possible $m$ is therefore
\[
\boxed{o(g^k)=\frac{n}{\gcd(n,k)}}.
\]
When $k=0$, the formula gives $1$, as it should.
\end{proof}

In $D_n$, $o(r)=n$ and every $sr^k$ has order $2$ (by $srs=r^{-1}$, $(sr^k)^2=sr^ksr^k=r^{-k}r^k=e$, and a reflection is not the identity). In $\mathfrak S_n$, for distinct $a_1,\ldots,a_k$ with $k\ge1$, the \emph{$k$-cycle} $(a_1\,a_2\,\ldots\,a_k)$ sends $a_1$ to $a_2$, each following entry to the next, and $a_k$ back to $a_1$, fixing all other elements. Its order is $k$: after $j$ applications, $a_1$ has advanced $j$ places, and first returns after $k$.

A single element generates a cyclic subgroup, and its powers eventually return to the identity when the group is finite. The next counting argument explains why that return time must divide the size of the whole group.

\section{Lagrange's theorem}

\begin{think}
In $\mathfrak S_3$, the subgroup $H=\langle(1\,2)\rangle=\{e,(1\,2)\}$ has $2$ elements. For each of the $6$ permutations $g$, multiply every element of $H$ on the left by $g$ and call the resulting set $gH$. How many different sets do you get, and how big is each? What does that suggest about $|H|$ and $|\mathfrak S_3|$?
\end{think}

\begin{definition}
Let $H\le G$. The \emph{left cosets} of $H$ are the sets $gH=\{gh:h\in H\}$, $g\in G$. Their number (possibly infinite) is the \emph{index} $[G:H]$.
\end{definition}

\begin{theorem}[Lagrange]\label{thm:lagrange}
Let $G$ be a finite group and $H\le G$. The left cosets of $H$ partition $G$, and each has $|H|$ elements. Hence
\[
\begin{gathered}|G|=[G:H]\cdot|H|,\qq\text{and in particular }|H|\text{ divides }|G|.\end{gathered}
\]
\end{theorem}
\begin{proof}
Fix an element $g\in G$. The map $h\mapsto gh$ from $H$ to the left coset $gH$ has an inverse: send $gh$ back to $g^{-1}(gh)=h$. Every left coset therefore contains exactly $|H|$ elements. We still have to show that we can count the cosets without counting any element twice.

Suppose $gH$ and $kH$ share an element $x$. Then $x=gh_1=kh_2$ for some $h_1,h_2\in H$. Thus $g=kh_2h_1^{-1}=kh$ for an element $h\in H$. For any $u\in H$, we have $gu=k(hu)\in kH$. Hence $gH\subseteq kH$. Interchanging $g$ and $k$ proves the reverse inclusion. Two cosets that meet are equal.

Every element $g$ belongs to $gH$, because $e\in H$. So the left cosets cover $G$, and distinct cosets are disjoint. If there are $[G:H]$ of them, each has $|H|$ elements. Counting the elements of $G$ gives
\[
|G|=\underbrace{|H|+\cdots+|H|}_{[G:H]\text{ cosets}}
=[G:H]|H|.
\]
In particular $|H|$ divides $|G|$.
\end{proof}

\paragraph{Why this works.}
The four rotations of a square form a subgroup $H$ of its eight symmetries.
If $s$ is a reflection, $sH$ consists of the other four symmetries; the
bijection $h\mapsto sh$ counts them. These two disjoint blocks account for
$8=2\cdot4$. Only the coset $H$ contains the identity, so the other coset
is not itself a subgroup. The theorem forbids a subgroup of order $3$,
but existence for a permitted divisor is a different question.

\begin{corollary}\label{cor:lagrange}
Let $G$ be finite of order $n$; in part (b), let $m\ge1$.
\begin{enumerate}[label=(\alph*)]
\item The order of every element divides $n$, and $g^n=e$ for all $g\in G$.
\item (Euler--Fermat, again) For $a$ coprime to $m$, $a^{\varphi(m)}\equiv1\pmod m$.
\item If $n$ is prime, $G$ is generated by any of its elements $g\ne e$.
\end{enumerate}
\end{corollary}
\begin{proof}
\emph{(a) Orders divide the size of the group.}
For any $g\in G$, two of $e,g,\ldots,g^n$ coincide, since there are $n+1$ terms in a set of size $n$. Thus $g$ has finite order by Proposition~\ref{prop:order}. Its powers form the subgroup
$\langle g\rangle$. Proposition~\ref{prop:order} shows
that this subgroup has exactly $o(g)$ elements.
Lagrange's theorem therefore gives $o(g)\mid n$.
Write $n=o(g)m$. Then
\[
g^n=(g^{o(g)})^m=e^m=e.
\]

\emph{(b) Euler's theorem as a group statement.}
Consider the group of invertible residue classes
$(\mathbb Z/m\mathbb Z)^*$, where $m\ge2$.
It has $\varphi(m)$ elements by the definition of
Euler's totient function. If $\gcd(a,m)=1$, the
class of $a$ belongs to this group. Apply (a):
\[
(a\bmod m)^{\varphi(m)}=1\bmod m,
\]
or $a^{\varphi(m)}\equiv1\pmod m$.
This is the Euler--Fermat theorem already proved in Volume~2,
now seen as a consequence of coset counting.
For $m=1$, every pair of integers is congruent modulo $1$, so the congruence holds there too.

\emph{(c) Prime-order groups.}
Suppose $|G|=n$ is prime and $g\ne e$.
The order $o(g)$ divides $n$ by (a) and is not $1$,
because an element of order $1$ is the identity.
The only remaining divisor is $n$.
Thus $\langle g\rangle$ contains all $n$ elements of $G$,
so $G=\langle g\rangle$.
\end{proof}

Euler's theorem in Volume~2 multiplied all the invertible residues together;
Lagrange now gives it for every finite group. We will see in Chapter~4 why
Lagrange's theorem has no converse: the permutation group $A_4$ has order
$12$ but no subgroup of order $6$.

Lagrange tells us what sizes are possible. For a group generated by one element, we can say more: every element is a power of the generator, and the entire multiplication table is determined by its order.

\section{Cyclic groups and isomorphisms}

Two groups can be ``the same'' while their elements have different names. The rotations $\{e,r,r^2,r^3\}$ of a square, the group $\mathbb U_4=\{1,i,-1,-i\}$ and $\Z/4\Z$ have the same multiplication table, up to renaming.

\begin{definition}
A map $f:G\to G'$ with $f(ab)=f(a)f(b)$ for all $a,b$ is a \emph{homomorphism}: it respects the group operation. A bijective homomorphism is an \emph{isomorphism}. Then $f(e)=e'$, $f(a^{-1})=f(a)^{-1}$, and $f^{-1}$ is an isomorphism too. If one exists, $G$ and $G'$ are \emph{isomorphic}, written $G\cong G'$. A group is \emph{cyclic} if it is generated by one element.
\end{definition}

Indeed $f(e)=f(ee)=f(e)f(e)$ gives $f(e)=e'$ by cancellation, then
$f(a)f(a^{-1})=f(e)=e'$. For a bijective $f$, the two sides of
$f^{-1}(xy)=f^{-1}(x)f^{-1}(y)$ have the same image $xy$, so injectivity
makes them equal. Isomorphic groups share properties defined from the
operation: order, being abelian or cyclic, element orders and subgroup counts.
For example, $f(a^m)=f(a)^m$ preserves element orders and generators;
$H\mapsto f(H)$ bijects the subgroups, with inverse $K\mapsto f^{-1}(K)$.

\begin{theorem}[cyclic groups]\label{thm:cyclic}
Let $G=\langle g\rangle$ be cyclic.
\begin{enumerate}[label=(\alph*)]
\item If $G$ is infinite, $k\mapsto g^k$ is an isomorphism $\Z\to G$. If $|G|=n$, $k\bmod n\mapsto g^k$ is an isomorphism $\Z/n\Z\to G$.
\item Every subgroup of $G$ is cyclic. If $|G|=n$, then for every divisor $d$ of $n$ there is exactly one subgroup of order $d$, namely $\langle g^{n/d}\rangle$, and there are no others.
\item If $|G|=n$, $g^k$ generates $G$ if and only if $\gcd(k,n)=1$. So $G$ has exactly $\varphi(n)$ generators, and exactly $\varphi(d)$ elements of order $d$ for each $d\mid n$.
\end{enumerate}
\end{theorem}
\begin{proof}
\emph{(a) Describe the group using integer exponents.}
Define $\Phi:\mathbb Z\to G$ by $\Phi(k)=g^k$.
The exponent rule gives
\[
\Phi(k+\ell)=g^{k+\ell}=g^kg^\ell
=\Phi(k)\Phi(\ell),
\]
so $\Phi$ is a homomorphism. It is onto by the
definition of $\langle g\rangle$.
If $G$ is infinite, its powers are all distinct by
Proposition~\ref{prop:order}; hence $\Phi$ is injective.
This proves $\mathbb Z\cong G$.

If $G$ has finite order $n$, the same proposition
shows that
\[
g^k=g^\ell\iff k\equiv\ell\pmod n.
\]
Thus the map $[k]\mapsto g^k$ on $\mathbb Z/n\mathbb Z$
is well defined, one-to-one and onto. It also
respects the operations, giving the claimed isomorphism.

\emph{(b) Every subgroup is cyclic.}
Let $H\le G$. The trivial subgroup is $\langle e\rangle$,
so suppose $H$ is not trivial. There is then a positive
integer $m$ with $g^m\in H$: if $H$ contains a non-identity
$g^k$ with negative exponent, take its inverse.
Choose the \emph{least} such positive $m$.

For any $g^k\in H$, divide the integer $k$ by $m$:
\[
k=qm+r,\qquad0\le r<m.
\]
Since $H$ contains $g^k$ and $g^m$, it contains
$g^r=g^k(g^m)^{-q}$. The minimality of $m$
forces $r=0$. Thus every element of $H$ is a power
of $g^m$, and conversely all such powers lie in $H$.
So $H=\langle g^m\rangle$.

When $|G|=n$, the identity $g^n=e\in H$ gives
$m\mid n$ by the same division argument. The order
formula of Proposition~\ref{prop:order} gives
\[
|H|=o(g^m)=\frac{n}{m}.
\]
A subgroup of order $d$ must therefore have
$m=n/d$ and is exactly $\langle g^{n/d}\rangle$.
Conversely, if $d\mid n$, the subgroup generated
by $g^{n/d}$ has order $d$. There is exactly
one subgroup of each divisor order.

\emph{(c) Count the generators.}
Again by the order formula,
\[
o(g^k)=\frac{n}{\gcd(n,k)}.
\]
This equals $n$ precisely when $\gcd(n,k)=1$.
Among the $n$ residue classes for $k$, there are
$\varphi(n)$ such classes, so $G$ has exactly
$\varphi(n)$ generators.

For any $d\mid n$, every element of order $d$
generates the unique subgroup of order $d$
from (b), and every generator of that subgroup
has order $d$. Applying the same reasoning to
that cyclic subgroup gives exactly $\varphi(d)$
elements of order $d$.

For example, $\mathbb Z/12\mathbb Z$ has a unique
subgroup of order $4$, generated by the class of $3$,
but that subgroup has only $\varphi(4)=2$ generators:
the classes of $3$ and $9$.
\end{proof}

\paragraph{Why this works.}
The exponent map $\Z\to\langle g\rangle$, $k\mapsto g^k$, is a
bijection when the powers never repeat. If they first return after $n$
steps, it instead gives a bijection from $\Z/n\Z$. For a subgroup $H$,
the allowed exponents $\{k:g^k\in H\}$ form an additive subgroup of
$\Z$, hence are $d\Z$ for some $d\ge0$ by the integer-subgroup
argument in the proof. The case $d=0$ means that only exponent $0$ is
allowed; it occurs for the trivial subgroup of an infinite cyclic group.
In a cyclic group of order $12$, the unique subgroup of order $4$ is
$\langle g^3\rangle$, with generators $g^3$ and $g^9$, since $\varphi(4)=2$.

In $\Z/n\Z$, for each $d\mid n$ the unique subgroup of order $d$
has $\varphi(d)$ generators, exactly the elements of order $d$.
These disjoint sets account for all $n$ elements, giving Gauss' identity
$\sum_{d\mid n}\varphi(d)=n$, proved in Volume~2 by another route.

\begin{beauty}
Once a group is cyclic of order $n$, the question ``which powers generate
it?'' becomes the integer test $\gcd(k,n)=1$. Thus $\varphi(n)$ counts
both invertible residues modulo $n$ and generators of every cyclic group
of order $n$, including rotations and roots of unity. For a finite field,
one step still needs proof: is its multiplicative group cyclic at all?
Exercise~\ref{ex:fieldcyclic} answers that using the polynomial root bound;
Chapter~9 will use the same argument to study finite fields.
\end{beauty}

\section{Exercises}

\begin{exercise}[$\star$]\label{ex:square}
Let $G$ be a group in which $x^2=e$ for every $x$. Show that $G$ is abelian.
\end{exercise}

\begin{exercise}[$\star$]\label{ex:dihedralcentre}
In $D_n$, show that the rotations form a subgroup of index $2$, and find the \emph{centre} $Z(D_n)=\{z:zg=gz\text{ for all }g\}$. (The answer depends on the parity of $n$.)
\end{exercise}

\begin{exercise}[$\star$]\label{ex:order4}
Show that a group of order $4$ is either cyclic or has all its non-identity elements of order $2$. Write the multiplication table in the second case, and give two geometric examples of it.
\end{exercise}

\begin{exercise}[$\star\star$]\label{ex:union}
Show that a group is never the union of two proper subgroups. Give a group that is the union of three proper subgroups.
\end{exercise}

\begin{exercise}[$\star\star$]\label{ex:gl2}
Let $p$ be prime. Show that $|GL_2(\Z/p\Z)|=(p^2-1)(p^2-p)$ and $|SL_2(\Z/p\Z)|=p(p^2-1)$. Check that $GL_2(\Z/2\Z)$ has order $6$ and is not abelian.
\end{exercise}

\begin{exercise}[$\star\star$]\label{ex:cubeS4}
\emph{The cube is $\mathfrak S_4$.} Number the four long diagonals of the cube $1,2,3,4$. Each rotation of the cube permutes the diagonals, which gives a map $\rho$ from the rotation group $C$ of Proposition~\ref{prop:cube} to $\mathfrak S_4$. Show that $\rho$ is an isomorphism.
\end{exercise}

\begin{exercise}[$\star\star$]\label{ex:prodorder}
Let $a,b$ be commuting elements of finite orders $m$ and $n$ with $\gcd(m,n)=1$. Show that $o(ab)=mn$. Show by an example in $\mathfrak S_3$ that the hypothesis $ab=ba$ cannot be dropped.
\end{exercise}

\begin{exercise}[$\star\star$]\label{ex:fieldcyclic}
Let $K$ be a field and $G$ a finite subgroup of $K^*$, of order $n$. Show that $G$ is cyclic. \emph{(For $K=\Z/p\Z$ this is the existence of primitive roots, Volume~2, Chapter~5.)}
\end{exercise}

\begin{exercise}[$\star$]\label{ex:wilson}
\emph{Wilson, from groups.} Let $G$ be a finite abelian group, written multiplicatively.
\begin{enumerate}[label=(\alph*)]
\item Show that the product of all the elements of $G$ equals the product of its elements of order $2$.
\item If $G$ has exactly one element $t$ of order $2$, show that this product is $t$; if it has none, the product is $e$.
\item Deduce Wilson's theorem: $(p-1)!\equiv-1\pmod p$ for every prime $p$.
\end{enumerate}
\end{exercise}

\begin{reminder}[polynomials and eigenvalues over a finite field]
The coordinate definitions from Volume~1 work over $K=\mathbb F_p$ too.
A scalar matrix is $\lambda I$. An eigenvalue $\lambda\in K$ has a
nonzero vector $v$ with $Mv=\lambda v$; equivalently,
$\det(\lambda I-M)=0$. For a $2\times2$ matrix the characteristic
polynomial is $\det(XI-M)=X^2-(\operatorname{tr}M)X+\det M$, where
$\operatorname{tr}M$ is the sum of the diagonal entries.
It is \emph{irreducible over $K$} if it has no factorization into two
polynomials of positive degree with coefficients in $K$. For this monic
quadratic that is equivalent to having no root in $K$: a factorization
must have two linear factors, while a root $r$ gives
\[
X^2-tX+d=(X-r)(X+r-t),\qquad d=tr-r^2.
\]
This uses no division by $2$, so it also works when $p=2$.
Uniformly choosing a matrix means giving each of its $p^4$ possible
entry lists the same mass $p^{-4}$.
\end{reminder}

\begin{exercise}[$\star\star$]\label{pb:matrices2}
Let $p$ be a prime, and let $M$ be a uniformly random matrix of $M_2(\mathbb F_p)$.
\begin{enumerate}[label=(\alph*)]
\item Count the matrices that are scalar; that have two distinct eigenvalues in $\mathbb F_p$; whose characteristic polynomial
is irreducible over $\mathbb F_p$; that are not scalar but have a double eigenvalue.
\item Deduce that, as $p\to\infty$, $M$ is diagonalisable over $\mathbb F_p$ with probability tending to $\frac12$, and has an
irreducible characteristic polynomial with probability tending to $\frac12$.
\end{enumerate}
\end{exercise}

\begin{exercise}[$\star$]\label{pb:fibdivides}
\emph{Every integer divides a Fibonacci number.}
\begin{enumerate}[label=(\alph*)]
\item Let $m\ge1$. Show that the sequence $(F_n\bmod m)$ is periodic from the start, and deduce that $m$ divides some $F_n$ with
$n\ge1$.
\item Show that the set of $n$ with $m\mid F_n$ is the set of multiples of the least such $n\ge1$.
\item Let $J_0=0$, $J_1=1$, $J_{n+2}=J_{n+1}+2J_n$. Show that $J_n$ is odd for every $n\ge1$. Which step of (a) fails?
\end{enumerate}
\end{exercise}

\begin{exercise}[$\star\star$]\label{pb:ruzsa}
\emph{Ruzsa's inequalities.} Let $A,B,C$ be finite sets in an abelian group, with $A\ne\varnothing$ for the covering statement in part (b). Write $A-B=\{a-b:a\in A,\,b\in B\}$ and $A+B=\{a+b:a\in A,\,b\in B\}$. These are sets of differences and sums, not the set-theoretic difference $A\setminus B$.
\begin{enumerate}[label=(\alph*)]
\item (Triangle inequality.) Show that $|A|\,|B-C|\le|A-B|\,|A-C|$. Deduce that if $|A+A|\le K|A|$ then $|A-A|\le K^2|A|$.
\item (Covering lemma.) Show that if $|A+B|\le K|A|$, then $B$ is covered by at most $K$ translates of $A-A$.
\end{enumerate}
\end{exercise}

\begin{exercise}[$\star\star\star$]\label{pb:returnbound}
\emph{The identity is the likeliest place.} Let $X$ be a symmetric random walk on a finite group $G$, started at the identity $e$,
and write $p_t(g)=P(X_t=g)$. This means $X_0=e$ and $X_t=\xi_1\cdots\xi_t$, where the increments are independent, all have the same distribution $\mu$, and $\mu(g)=\mu(g^{-1})$. Times $t$ are nonnegative integers; all probabilities here concern finite experiments.
\begin{enumerate}[label=(\alph*)]
\item Show that $p_{2t}(e)=\sum_{g}p_t(g)^2$, and deduce that $p_{2t}(g)\le p_{2t}(e)$ for every $g\in G$ and that $p_{2t}(e)\ge\frac1{|G|}$.
\item Show that $t\mapsto p_{2t}(e)$ is non-increasing.
\item Show that $p_{2t+1}(e)$ can be $0$ for every $t$, but that $p_{2t+1}(g)\le p_{2t}(e)$ for all $g$ and $t$.
\item For $n\ge1$, the walk on $\Z/n\Z$ with equally likely steps $\pm1$, started at $0$, is at time $2t$ at $0$ with probability $\frac1n\sum_{k=0}^{n-1}\cos^{2t}\frac{2\pi k}n$. Derive this from the finite geometric sum for roots of unity, then check (a) and
(b) directly on this formula.
\end{enumerate}
\end{exercise}

\section{Hints}

\begin{hint}{ex:square}
$x^2=e$ means $x=x^{-1}$. Apply it to $x$, $y$ and $xy$.
\end{hint}
\begin{hint}{ex:dihedralcentre}
Use $sr^k=r^{-k}s$. When does $r^k$ commute with $s$? Can a reflection commute with $r$?
\end{hint}
\begin{hint}{ex:order4}
By Lagrange, the orders of elements are $1$, $2$ or $4$. Without an element of order $4$, use Proposition~\ref{prop:rules}(d) to fill the table. For examples, think of a rectangle, or of $(\Z/8\Z)^*$.
\end{hint}
\begin{hint}{ex:union}
If $G=H\cup K$ with $h\in H\setminus K$ and $k\in K\setminus H$, where is $hk$? For three subgroups, try a group of order $4$.
\end{hint}
\begin{hint}{ex:gl2}
Choose the first column (any non-zero vector), then the second (any vector not on the line of the first). For $SL_2$, the determinant is onto $(\Z/p\Z)^*$, and each value is taken equally often.
\end{hint}
\begin{hint}{ex:cubeS4}
$\rho$ respects products because both sides are composition. Since both groups have $24$ elements, it is enough to show that only the identity fixes every diagonal. Use the signed permutation matrices of the proof of Proposition~\ref{prop:cube}.
\end{hint}
\begin{hint}{ex:prodorder}
$(ab)^k=a^kb^k$. If $(ab)^k=e$, then $a^k=b^{-k}$ lies in $\langle a\rangle\cap\langle b\rangle$. What can the order of an element of this intersection be?
\end{hint}
\begin{hint}{ex:fieldcyclic}
Let $\psi(d)$ be the number of elements of order $d$ in $G$. If $\psi(d)\ne0$, an element of order $d$ generates $d$ distinct roots of $X^d-1$, which has at most $d$ roots in a field. Compare $\psi(d)$ with $\varphi(d)$ and use $\sum_{d\mid n}\varphi(d)=n$.
\end{hint}
\begin{hint}{ex:wilson}
(a) Pair each element with its inverse. (b) The elements of order $2$, together with $e$, form a subgroup. (c) In $(\Z/p\Z)^*$, $x^2=1$ means $p\mid(x-1)(x+1)$.
\end{hint}

\begin{hint}{pb:matrices2}
Group matrices by a change of basis, $C\mapsto PCP^{-1}$. Two changes of basis give the same matrix exactly when their ratio commutes with $C$; use Lagrange to count the equal fibres. For a non-scalar $2\times2$ matrix, choose a basis $(v,Cv)$ and show that every commuting matrix is $aI+bC$.
Alternatively, fix the trace $t$ and determinant $d$, write the diagonal
entries as $a,t-a$, and count the solutions of $bc=a(t-a)-d$.
The answer changes according to whether $a$ is a root of $X^2-tX+d$.
\end{hint}

\begin{hint}{pb:fibdivides}
(a) The map $(u,v)\mapsto(v,u+v)$ on $(\Z/m\Z)^2$ is a bijection. (b) $\gcd(F_a,F_b)=F_{\gcd(a,b)}$.
\end{hint}

\begin{hint}{pb:ruzsa}
(a) For each $x\in B-C$ fix one way of writing it. (b) Take a maximal set of disjoint translates $b+A$, $b\in B$.
\end{hint}

\begin{hint}{pb:returnbound}
(a) $X_{2t}=X_t\cdot(X_t^{-1}X_{2t})$, and the second factor has the law of $X_t$, inverted by symmetry. Then Cauchy--Schwarz, and $\sum_gp_t(g)=1$. (b) $p_{t+1}$ is
an average of translates of $p_t$, and the $\ell^2$ norm is convex. (c) Split $2t+1=t+(t+1)$.
\end{hint}

\section{Solutions}

\begin{solution}{ex:square}
\idea To show that a group is abelian, we need $xy=yx$ for any $x,y$. The hypothesis applies not only to $x$ and $y$ separately, but also to their product $xy$. Combine it with the rule that taking the inverse reverses multiplication order.

\writeup Let $x,y\in G$. The condition $x^2=e$ implies $x^{-1}=x$, and likewise $y^{-1}=y$. The condition also applies to the group element $xy$, giving
\[
(xy)^2=e\qquad\Longrightarrow\qquad(xy)^{-1}=xy.
\]
On the other hand, the inverse-of-a-product formula, valid in every group, gives
\[
(xy)^{-1}=y^{-1}x^{-1}=yx.
\]
Comparing the two expressions yields $xy=yx$. Because $x,y$ were arbitrary, every pair of elements commutes and $G$ is abelian.

The hypothesis must hold for \emph{every} group element, including products. An element of order $2$ is called an \emph{involution}. Reflections of a regular polygon are involutions, but two reflections need not commute; knowing the hypothesis only for some chosen elements would not suffice.
\end{solution}

\begin{solution}{ex:dihedralcentre}
\idea Every element of $D_n$ is either a rotation $r^k$ or a reflection $sr^k$. A central element must commute with the two generators $r$ and $s$; there is no need to compare it with all $2n$ elements separately.

\writeup Recall that $D_n$ is generated by $r,s$, with $r^n=s^2=e$ and $srs=r^{-1}$, where $n\ge3$.
The $n$ rotations $e,r,\ldots,r^{n-1}$ form the subgroup $\langle r\rangle$, and the $n$ reflections are the elements $sr^k$. Thus $|D_n|=2n$ and
\[
[D_n:\langle r\rangle]=\frac{2n}{n}=2.
\]

Let us now find the centre. The relation $sr=r^{-1}s$ implies, by induction, $sr^k=r^{-k}s$ for every integer $k$.

A rotation $r^k$ automatically commutes with $r$. It commutes with $s$ exactly when
\[
sr^k=r^ks
\quad\Longleftrightarrow\quad
r^{-k}s=r^ks
\quad\Longleftrightarrow\quad
r^{2k}=e
\quad\Longleftrightarrow\quad n\mid2k.
\]
If $n$ is odd, the only residue $k$ satisfying this is $k=0$. If $n$ is even, there are two, $k=0$ and $k=n/2$.

Could a reflection be central? Take $sr^k$ and compare its products with $r$:
\[
(sr^k)r=sr^{k+1},
\qquad
r(sr^k)=sr^{k-1}.
\]
Equality would give $r^2=e$, impossible because $r$ has order $n\ge3$. No reflection is central.

We have therefore exhausted the elements of $D_n$:
\[
\boxed{Z(D_n)=
\begin{cases}
\{e\},&n\text{ odd},\\
\{e,r^{n/2}\},&n\text{ even}.
\end{cases}}
\]
Geometrically, in an even-sided regular polygon the half-turn commutes with every symmetry; in an odd-sided one there is no non-trivial central rotation.
\end{solution}

\begin{solution}{ex:order4}
\idea Start with Lagrange: an element can have order $1$, $2$, or $4$. If nobody has order $4$, label the three non-identity elements $a,b,c$ and let cancellation determine every product. For the examples, look for geometric transformations that are their own inverses.

\writeup Let $|G|=4$. If some $g\in G$ has order $4$, then $\langle g\rangle$ already has all four elements, and $G\cong\Z/4\Z$.

Otherwise the order of each non-identity element is $2$, so write
\[
G=\{e,a,b,c\},\qquad a^2=b^2=c^2=e.
\]
The product $ab$ cannot equal $e$: that would give $b=a^{-1}=a$. It cannot equal $a$ or $b$, because cancellation would make $b=e$ or $a=e$. The only remaining element is $c$, so $ab=c$. The same argument, with $a$ and $b$ interchanged, shows $ba=c$. Repeating it for the other pairs forces the entire table:
\[
\begin{array}{c|cccc}
 &e&a&b&c\\\hline
e&e&a&b&c\\
a&a&e&c&b\\
b&b&c&e&a\\
c&c&b&a&e
\end{array}.
\]
This group is the \emph{Klein four-group}, often written $V_4$ or $(\Z/2\Z)^2$. Every row can be checked directly from the identities above.

Here are two \emph{geometric} copies. A non-square rectangle in the plane has exactly four symmetries: the identity, reflection in its horizontal midline, reflection in its vertical midline, and the half-turn. Each non-identity symmetry has order $2$. In three dimensions, the identity and the half-turns around the $x$-, $y$-, and $z$-axes form another $V_4$: composing the half-turns about two perpendicular coordinate axes gives the half-turn about the third. Both examples have precisely the displayed table.
\end{solution}

\begin{solution}{ex:union}
\idea A union of two subgroups fails closure unless one subgroup contains the other. Choose an element from each subgroup that does not belong to the other, then look at their product.

\writeup Suppose, toward a contradiction, that $G=H\cup K$ with $H$ and $K$ both proper subgroups.

If $H\subseteq K$, then $G=K$, contradicting that $K$ is proper. The same argument excludes $K\subseteq H$. Consequently there exist
\[
h\in H\setminus K,\qquad k\in K\setminus H.
\]
Since $hk\in G=H\cup K$, either $hk\in H$ or $hk\in K$.

If $hk\in H$, closure under multiplication and inverses gives
\[
k=h^{-1}(hk)\in H,
\]
contrary to our choice of $k$. If $hk\in K$, the analogous calculation
\[
h=(hk)k^{-1}\in K
\]
contradicts our choice of $h$. Neither case is possible. Therefore two proper subgroups never cover a group.

Three proper subgroups \emph{can} cover one. Let $V_4=\{e,a,b,c\}$ be the Klein four-group of Exercise~\ref{ex:order4}, where every non-identity element has order $2$. Then
\[
H_a=\{e,a\},\qquad H_b=\{e,b\},\qquad H_c=\{e,c\}
\]
are each subgroups of order $2$, hence proper, and
\[
V_4=H_a\cup H_b\cup H_c.
\]
They overlap at the identity. This is precisely what allowed three subgroups to succeed where two could not.
\end{solution}

\begin{solution}{ex:gl2}
\idea A $2\times2$ matrix is invertible if and only if its two column vectors form a basis of the two-dimensional vector space over $\mathbb F_p$. Choose the first vector, then the second. For determinant $1$, use that all non-zero determinant values occur equally often.

\writeup Put $K=\mathbb F_p=\Z/p\Z$. The first column of an invertible matrix can be any non-zero vector of $K^2$, so there are $p^2-1$ choices. Once a first column $v\ne0$ is chosen, the second column must lie outside the line $\langle v\rangle$, which has exactly $p$ vectors. Hence there are $p^2-p$ choices for the second column. Every choice gives a different invertible matrix, so
\[
\boxed{|GL_2(K)|=(p^2-1)(p^2-p)}.
\]

For every $\lambda\in K^*$, left multiplication by $\operatorname{diag}(\lambda,1)$ is a bijection from the matrices of determinant $1$ to those of determinant $\lambda$. Its inverse is multiplication by $\operatorname{diag}(\lambda^{-1},1)$. The $p-1$ possible nonzero determinants therefore occur equally often. Partitioning $GL_2(K)$ by determinant gives
\[
|SL_2(K)|=\frac{|GL_2(K)|}{p-1}
=\frac{(p^2-1)p(p-1)}{p-1}
=\boxed{p(p^2-1)}.
\]

For $p=2$, the formula gives $(4-1)(4-2)=3\cdot2=6$ invertible matrices. This group is not abelian. Indeed, in $\mathbb F_2$ let
\[
A=\begin{pmatrix}1&1\\0&1\end{pmatrix},
\qquad
B=\begin{pmatrix}1&0\\1&1\end{pmatrix}.
\]
Both have determinant $1$, yet
\[
AB=\begin{pmatrix}0&1\\1&1\end{pmatrix},
\qquad
BA=\begin{pmatrix}1&1\\1&0\end{pmatrix}.
\]
These products differ, proving non-commutativity. We will later recognise this six-element group as another copy of $\mathfrak S_3$.
\end{solution}

\begin{solution}{ex:cubeS4}
\idea An orientation-preserving symmetry of the cube permutes its four long diagonals. That gives a homomorphism to $\mathfrak S_4$. Both groups have $24$ elements, so it is enough to prove injectivity. For this, show that a rotation fixing every diagonal \emph{as a line} must be the identity.

\writeup Centre the cube at the origin with vertices $(\pm1,\pm1,\pm1)$. Its four long diagonals join opposite vertices. A rotation carries opposite vertices to opposite vertices and hence permutes those four lines. Composition of rotations induces composition of line permutations, giving
\[
\rho:C\longrightarrow\mathfrak S_4.
\]

We know from Proposition~\ref{prop:cube} that $|C|=24$. Since $|\mathfrak S_4|=24$ too, it is enough to show that only the identity rotation fixes every diagonal. Indeed, if $\rho(U)=\rho(V)$, the rotation $V^{-1}U$ fixes every diagonal; if that forces $V^{-1}U=I$, then $U=V$.

Let $T$ fix each long diagonal as an unoriented line. By Proposition~\ref{prop:cube}, every cube rotation is a signed permutation matrix of determinant $+1$. Take the diagonal through $v=(1,1,1)$. Fixing it as a line means
\[
T(1,1,1)=\varepsilon(1,1,1)
\]
for some $\varepsilon=\pm1$. In a signed permutation matrix, $T(1,1,1)$ simply lists the three signs of its non-zero entries. Hence all three signs are $\varepsilon$, and
\[
T=\varepsilon P
\]
for an ordinary permutation matrix $P$.

Next consider the diagonal through $w=(1,1,-1)$. The vector $Pw$ has exactly one negative coordinate. Because $T$ preserves the line of $w$, we have $Pw=\pm w$. But $-w=(-1,-1,1)$ has \emph{two} negative coordinates, so $Pw$ cannot equal $-w$. Therefore $Pw=w$ and $P$ leaves the third coordinate fixed. Apply the same argument to $(1,-1,1)$: $P$ leaves the second coordinate fixed. Hence it also fixes the first, and $P=I$.

We have proved $T=\varepsilon I$. Because $T$ is a rotation,
\[
1=\det T=\varepsilon^3,
\]
forcing $\varepsilon=1$. Thus $T=I$, so the kernel is trivial.

The preceding injectivity test now applies to $\rho$. An injection between finite sets of the same size $24$ is bijective; hence
\[
\boxed{C\cong\mathfrak S_4}.
\]
The isomorphism is geometric: permuting four diagonals packages every spatial cube rotation into a permutation of just four objects.
\end{solution}

\begin{solution}{ex:prodorder}
\idea If $a$ and $b$ commute, powers of their product split as powers of the factors. If $(ab)^k=e$, then $a^k=b^{-k}$ is trapped in two cyclic subgroups whose orders are coprime. That common element can only be the identity.

\writeup Suppose $ab=ba$, $o(a)=m$, $o(b)=n$, and $\gcd(m,n)=1$. The commutation relation implies, by induction on positive $k$,
\[
(ab)^k=a^kb^k.
\]
In particular,
\[
(ab)^{mn}=(a^m)^n(b^n)^m=e,
\]
so the order of $ab$ divides $mn$.

For the reverse divisibility, let $k\ge1$ satisfy $(ab)^k=e$. Then
\[
a^k=b^{-k}.
\]
Call this common value $x$. It lies in both $\langle a\rangle$ and $\langle b\rangle$, so Lagrange's theorem applied to those subgroups says that $o(x)$ divides both $m$ and $n$. As the two integers are coprime, $o(x)=1$, hence $x=e$. Thus $a^k=b^k=e$, which is equivalent to
\[
m\mid k,\qquad n\mid k.
\]
Coprimality then implies $mn\mid k$. Taking $k=o(ab)$ proves $mn\mid o(ab)$, while the first paragraph gave the reverse divisibility. Therefore
\[
\boxed{o(ab)=mn}.
\]

The commuting hypothesis cannot be discarded. In $\mathfrak S_3$ put $a=(1\,2)$ of order $2$ and $b=(1\,2\,3)$ of order $3$. With composition from right to left,
\[
ab=(1\,2)(1\,2\,3)=(2\,3),
\]
which has order $2$, not $6$. Indeed $a$ and $b$ do not commute. Coprime orders alone do not give the product formula.
\end{solution}

\begin{solution}{ex:fieldcyclic}
\idea A field polynomial has no more distinct roots than its degree. That simple fact sharply restricts the number of group elements of each order. Count the elements of $G$ by their orders and compare with the familiar divisor-sum identity for Euler's totient.

\writeup Let $G$ be a finite subgroup of $K^*$, with $|G|=n$. Since multiplication in a field is commutative, $G$ is abelian. For each divisor $d$ of $n$, let
\[
\psi(d)=|\{g\in G:o(g)=d\}|.
\]
By Lagrange's theorem, every element has order dividing $n$, so the sets counted by $\psi(d)$ partition $G$:
\[
\sum_{d\mid n}\psi(d)=n.
\]

We claim that for each $d\mid n$, either $\psi(d)=0$ or $\psi(d)=\varphi(d)$. If no element has order $d$, the first alternative holds. Otherwise choose $g\in G$ of order $d$. The elements
\[
1,g,g^2,\ldots,g^{d-1}
\]
are $d$ \emph{distinct} roots of $X^d-1$. The polynomial root bound holds over this field too. Indeed, $F(X)-F(a)$ is divisible by $X-a$, by the identity for $X^j-a^j$. If $F(a)=0$, write $F=(X-a)Q$; every other root is a root of $Q$, because its difference from $a$ is nonzero and can be cancelled. Induction on the degree proves that a nonzero degree-$d$ polynomial has at most $d$ distinct roots. Thus the listed powers are \emph{all} the roots of $X^d-1$ in $K$. Any other element $h\in G$ of order $d$ satisfies $h^d=1$, so it lies in $\langle g\rangle$. In a cyclic group of order $d$, exactly $\varphi(d)$ elements have order $d$ (Theorem~\ref{thm:cyclic}). This proves the claim.

Consequently $\psi(d)\le\varphi(d)$ for each divisor of $n$. Summing these inequalities gives
\[
n=\sum_{d\mid n}\psi(d)
\le\sum_{d\mid n}\varphi(d)=n,
\]
where the last identity was proved in Volume~2. Since the two sums are equal and all terms are non-negative, \emph{every} individual inequality must be an equality:
\[
\psi(d)=\varphi(d)\quad(d\mid n).
\]

In particular, $\psi(n)=\varphi(n)\ge1$, so some $g\in G$ has order exactly $n$. Then $|\langle g\rangle|=n=|G|$, which forces
\[
\boxed{G=\langle g\rangle}.
\]
Thus every finite subgroup of the multiplicative group of a field is cyclic. For $K=\mathbb F_p$, this gives the existence of a primitive root modulo $p$.
\end{solution}

\begin{solution}{ex:wilson}
\idea When a group is abelian, we may multiply its elements in any order. Pair each element $g$ with its inverse $g^{-1}$; pairs of distinct elements disappear. Only the self-inverse elements remain, and they have order $1$ or $2$.

\writeup \emph{(a) Pairing inverse elements.} Since $G$ is finite and abelian, the product
\[
P=\prod_{g\in G}g
\]
does not depend on the order of its factors. Applying $g\mapsto g^{-1}$ twice returns $g$, so it partitions $G$ into singleton sets and pairs $\{g,g^{-1}\}$ of distinct elements. Each pair contributes $gg^{-1}=e$. The singleton sets are precisely those satisfying $g=g^{-1}$, or $g^2=e$.

Such elements are either $e$ (order $1$) or elements of order $2$. As $e$ does not affect the product, we have shown
\[
\boxed{P=\prod_{\substack{g\in G\\o(g)=2}}g},
\]
with an empty product interpreted as $e$.

\emph{(b) Zero or one involution.} If $t$ is the unique element of order $2$, the remaining product has one factor and equals $t$. If no element has order $2$, all factors cancel in inverse pairs and $P=e$. More generally, the self-inverse elements form a subgroup $T=\{g:g^2=e\}$ because $G$ is abelian and $(gh)^2=g^2h^2=e$; the two special cases just describe $|T|=1$ or $2$.

\emph{(c) Wilson's congruence.} Take the multiplicative group
\[
G=(\Z/p\Z)^*
\]
for a prime $p$. This is abelian, and its elements are the classes $1,2,\ldots,p-1$. The self-inverse elements satisfy
\[
x^2=1\quad\Longleftrightarrow\quad(x-1)(x+1)=0\quad\text{in }\mathbb F_p.
\]
Because a field has no zero divisors, $x=1$ or $x=-1$. If $p$ is odd, these are different and $-1$ is the unique element of order $2$. Part (b) therefore gives
\[
1\cdot2\cdots(p-1)=-1\quad\text{in }\mathbb F_p,
\]
that is,
\[
\boxed{(p-1)!\equiv-1\pmod p}.
\]
For $p=2$, the two classes $1$ and $-1$ coincide, but $1!=1\equiv-1\pmod2$, so the result still holds.

The proof shows Wilson's theorem as a consequence of finite-group pairing, rather than as a special factorial trick.
\end{solution}

\begin{solution}{pb:matrices2}
\idea Count how many changes of basis give each matrix; those that leave
$C$ unchanged commute with $C$. After that structural count, check the
answer by fixing the trace and determinant and counting the remaining entries.

\writeup Let $K=\mathbb F_p$ and $G=GL_2(K)$. Choosing two independent columns gives
\[
|G|=(p^2-1)(p^2-p)=p(p-1)^2(p+1).
\]

\emph{Why a non-scalar matrix has a simple centraliser.}
Let $C\in M_2(K)$ be non-scalar. There is a vector $v$ for which $v,Cv$ are linearly independent: otherwise every nonzero vector would be an eigenvector. For independent $u,w$, write $Cu=\alpha u$, $Cw=\beta w$ and $C(u+w)=\gamma(u+w)$. Comparing the independent coefficients gives $\alpha=\beta=\gamma$, which would make $C$ scalar. Thus $(v,Cv)$ is a basis. If $AC=CA$ and $Av=av+bCv$, then
\[
(A-aI-bC)v=0.
\]
The matrix $A-aI-bC$ commutes with $C$, so it also kills $Cv$ and hence both vectors of a basis. Therefore
\[
A=aI+bC.
\]
The matrices $aI+bC$ all commute with $C$, so this list is exact. Its $p^2$ members are distinct, since $aI+bC=0$ with $b\ne0$ would make $C$ scalar. The \emph{centraliser} in $G$ means the invertible members of this list. We now justify both the similarity count and its size in each case.

Before counting the kinds of matrices, justify the change-of-basis count. The invertible matrices commuting with $C$ form a subgroup $H$ of $G$: products and inverses still commute with $C$. For $P,Q\in G$,
\[
PCP^{-1}=QCQ^{-1}\iff Q^{-1}P\in H\iff P\in QH.
\]
Thus each resulting matrix is obtained from exactly $|H|$ changes of basis. The number of distinct matrices similar to $C$ is $|G|/|H|$, by the coset argument already proved in Lagrange's theorem. This derives the needed count without assuming the later theory of group actions.

Write the characteristic polynomial as $q(X)=X^2-tX+d$. Direct multiplication of a $2\times2$ matrix gives $C^2-tC+dI=0$. In the basis $(v,Cv)$ its columns are consequently $(0,1)^{\mathsf T}$ and $(-d,t)^{\mathsf T}$. This proves the assertion above that all non-scalar matrices with the same $q$ are similar. We can count the invertible matrices $aI+bC$ directly, without constructing a new field or a quotient ring.

\emph{Four disjoint kinds of matrices.}
There are $p$ scalars $\lambda I$, one for each $\lambda\in K$.

If $q$ has distinct roots $\lambda,\mu\in K$, choose nonzero eigenvectors for them. They are independent: if one were a nonzero multiple of the other, applying $C$ would force $\lambda=\mu$. In that basis $aI+bC$ is diagonal, with entries $a+b\lambda$ and $a+b\mu$. The map
\[
(a,b)\longmapsto(a+b\lambda,a+b\mu)
\]
is a bijection of $K^2$, since its determinant is $\mu-\lambda\ne0$. Exactly $(p-1)^2$ pairs give two nonzero entries. Hence each unordered choice of two distinct eigenvalues gives $|G|/(p-1)^2=p(p+1)$ matrices. There are $p(p-1)/2$ choices, for a total of $\frac12p^2(p^2-1)$.

If $q$ is irreducible, then
\[
\det(aI+bC)=a^2+tab+db^2.
\]
For $b\ne0$, this is $b^2q(-a/b)$, which cannot vanish because $q$ has no root in $K$. For $b=0$, it is $a^2$. Thus all $p^2-1$ nonzero pairs $(a,b)$ give invertible matrices, so each $q$ gives $|G|/(p^2-1)=p(p-1)$ similar matrices. There are $p^2$ monic quadratics in total. The reducible ones correspond to unordered pairs of roots with repetition, of which there are $p(p+1)/2$. The number of irreducible quadratics is therefore $(p^2-p)/2$, giving $\frac12p^2(p-1)^2$ matrices in this case.

Finally, if $q=(X-\lambda)^2$ and $C$ is non-scalar, put $N=C-\lambda I$. Then $N\ne0$ and $N^2=0$. The commuting matrices are
\[
aI+bC=(a+b\lambda)I+bN,
\]
with determinant $(a+b\lambda)^2$. There are $p(p-1)$ invertible choices. Each $\lambda$ therefore gives $|G|/(p(p-1))=p^2-1$ matrices, and the total is $p(p^2-1)$.

We obtain the requested count:
\[
\begin{array}{l|r}
\text{matrix type}&\text{number}\\\hline
\text{scalar}&p\\
\text{two distinct eigenvalues in }K&\frac12p^2(p^2-1)\\
\text{irreducible characteristic polynomial}&\frac12p^2(p-1)^2\\
\text{non-scalar repeated eigenvalue}&p(p^2-1).
\end{array}
\]
As a check, the four counts add up to $p^4$, the total number of $2\times2$ matrices over $K$.

\emph{A second count, directly from the entries.}
Fix $q(X)=X^2-tX+d$, and let $r_q\in\{0,1,2\}$ be its number of
distinct roots in $K$. A matrix with this characteristic polynomial is
\[
\begin{pmatrix}a&b\\c&t-a\end{pmatrix},\qquad
bc=a(t-a)-d=-q(a).
\]
If $q(a)\ne0$, each of the $p-1$ nonzero choices of $b$ gives exactly
one $c$. If $q(a)=0$, the pairs with $bc=0$ are the two coordinate axes,
with their common pair $(0,0)$ counted once: $2p-1$ pairs. Summing over $a$
gives
\[
(p-r_q)(p-1)+r_q(2p-1)=p(p-1)+pr_q.
\]
Thus a fixed irreducible quadratic gives $p(p-1)$ matrices, one with two
distinct roots gives $p(p+1)$, and $(X-\lambda)^2$ gives $p^2$.
The last count contains exactly one scalar, $\lambda I$, leaving
$p^2-1$ non-scalars. Multiply by the numbers of each kind of quadratic
already counted above, and add the $p$ scalars: this reproduces every entry
of the table without changes of basis. The only field division in this
direct count is by a nonzero $b$, so it also works for $p=2$; the halves
in the table are ordinary integer counts, not division by $2$ in $K$.

\emph{(b) The limiting probabilities.}
A matrix is diagonalisable over $K$ precisely when it is scalar or has two distinct eigenvalues in $K$. A non-scalar matrix with a repeated eigenvalue cannot be diagonalisable, because a diagonalisable matrix with just that eigenvalue would equal $\lambda I$. Consequently
\[
\Pr(M\text{ diagonalisable over }K)
=\frac{p+\frac12p^2(p^2-1)}{p^4}
\longrightarrow\frac12.
\]
And
\[
\Pr(\chi_M\text{ irreducible})
=\frac{\frac12p^2(p-1)^2}{p^4}
=\frac{(p-1)^2}{2p^2}\longrightarrow\frac12.
\]

\tools Changes of basis, the $2\times2$ determinant, and the equal-coset count in Lagrange's theorem. All counts needed here were derived above.

\leads The two permutations of two letters have cycle types $1+1$ and $2$.
The two limits $1/2$ match their equal probabilities. Does a similar relation
hold for higher-degree polynomials and permutations of more letters?
\end{solution}

\begin{solution}{pb:fibdivides}
\idea The number $F_n$ alone does not tell us the next Fibonacci number. The pair $(F_n,F_{n+1})$ does. Modulo $m$ that pair has only finitely many states, and its update rule can be reversed, which forces a cycle from the very beginning.

\writeup \emph{(a) Return to the starting pair.}
Fix $m\ge1$ and consider the map on $(\mathbb Z/m\mathbb Z)^2$,
\[
T(u,v)=(v,u+v).
\]
Given $(a,b)$ we can recover the previous state: it was $(b-a,a)$. Hence $T$ is bijective, with
\[
T^{-1}(a,b)=(b-a,a).
\]
Starting at $(F_0,F_1)=(0,1)$, successive applications of $T$ give exactly the consecutive Fibonacci pairs modulo $m$.

A bijection of a finite set is a permutation. Every orbit of a permutation is a cycle, including its starting point; it cannot first wander through non-repeating states and then enter a different cycle. Therefore some $P\ge1$ satisfies
\[
(F_P,F_{P+1})\equiv(0,1)\pmod m.
\]
In particular $F_P\equiv0\pmod m$, proving that every positive integer divides some positive-index Fibonacci number. When $m=1$ this is immediate, and the state argument still applies.

\emph{(b) The first divisible index controls all the others.}
Let $z=z(m)$ be the smallest positive index with $m\mid F_z$, which exists by (a). Volume~1 proved
\[
\gcd(F_a,F_b)=F_{\gcd(a,b)}
\qquad(a,b\ge0,\text{ not both zero}).
\]
Suppose $n\ge1$ and $m\mid F_n$. Since $m$ also divides $F_z$, it divides their greatest common divisor:
\[
m\mid F_{\gcd(n,z)}.
\]
The positive index $\gcd(n,z)$ is at most $z$, so the minimality of $z$ forces $\gcd(n,z)=z$, or $z\mid n$.

Conversely, if $z\mid n$, the same Fibonacci gcd identity yields
\[
\gcd(F_z,F_n)=F_z,
\]
which means $F_z\mid F_n$. Since $m\mid F_z$, also $m\mid F_n$. We have proved
\[
\boxed{m\mid F_n\iff z(m)\mid n.}
\]
The index $n=0$ satisfies both conditions too, because $F_0=0$ and every positive integer divides $0$.

\emph{(c) Why reversibility mattered.}
For $J_0=0$, $J_1=1$ and $J_{n+2}=J_{n+1}+2J_n$, reduction modulo $2$ gives
\[
J_{n+2}\equiv J_{n+1}\pmod2.
\]
As $J_1=1$ and $J_2=1$, induction shows $J_n\equiv1\pmod2$ for every $n\ge1$. In particular no positive-index term is divisible by $2$.

The corresponding state update is
\[
S(u,v)=(v,v+2u).
\]
Modulo $2$ it reduces to $S(u,v)=(v,v)$, which is not injective: $S(0,0)=S(1,0)$. So $S$ is not a permutation of the finite state set. Its orbit may enter a cycle without ever returning to its initial pair. That is exactly the step in (a) that no longer works.

\tools Fibonacci gcd and divisibility identities (Volume~1, Chapter~1); permutations and finite-state cycles.

\leads Whenever a linear recurrence has an invertible state-transition matrix modulo $m$, its state sequence is periodic from the start. The rank of apparition $z(m)$ captures the resulting divisibility pattern for Fibonacci numbers.
\end{solution}

\begin{solution}{pb:ruzsa}
\idea We want to compare three sets of differences. A well-chosen injection proves the inequality without estimating each set separately. For the covering statement, choose a maximal family of disjoint translates and see what maximality forces.

\writeup \emph{(a) The triangle inequality.}
If $B-C$ is empty, both sides of the claimed inequality are non-negative and the left side is zero, so there is nothing to prove. Otherwise, for each $x\in B-C$, choose once and for all one pair
\[
b_x\in B,\quad c_x\in C,\qquad x=b_x-c_x.
\]
Define
\[
\Phi:A\times(B-C)\longrightarrow(A-B)\times(A-C),
\quad
\Phi(a,x)=(a-b_x,\ a-c_x).
\]
The output belongs to the required product. We must prove that two distinct inputs never give the same output. If the output is $(u,v)$, subtract its coordinates:
\[
v-u=(a-c_x)-(a-b_x)=b_x-c_x=x.
\]
So the output determines $x$. Because our representation $(b_x,c_x)$ was fixed for each $x$, it then determines $a=u+b_x$. Hence $\Phi$ is injective, and counting its domain and codomain proves
\[
\boxed{|A|\,|B-C|\le|A-B|\,|A-C|.}
\]

Suppose now that $A$ is non-empty and $|A+A|\le K|A|$. Apply the inequality with $B=C=-A$. We get
\[
|A|\,|A-A|\le|A+A|^2\le K^2|A|^2.
\]
Divide by the positive integer $|A|$:
\[
\boxed{|A-A|\le K^2|A|.}
\]
If $A$ is empty, the same inequality is simply $0\le0$. Thus no separate assumption is needed for this consequence.

\emph{(b) The covering lemma.}
Assume $A\ne\varnothing$ and $|A+B|\le K|A|$. If $B$ is empty, the empty family of translates covers it. Otherwise choose a subset $X\subseteq B$ maximal with the property that the translates $x+A$, for $x\in X$, are pairwise disjoint. Such a maximal choice exists because $B$ is finite; for example, add translates one at a time until none can be added.

Every $x+A$ lies in $A+B$, and each has $|A|$ elements. Disjointness yields
\[
|X|\,|A|\le|A+B|\le K|A|.
\]
Since $|A|>0$, this implies $|X|\le K$.

Now take any $b\in B$. If $b+A$ were disjoint from every $x+A$ with $x\in X$, we could add $b$ to $X$, contradicting maximality. Thus there are $x\in X$ and $a_1,a_2\in A$ such that
\[
b+a_1=x+a_2.
\]
Rearranging in the abelian group gives
\[
b=x+(a_2-a_1)\in x+(A-A).
\]
This holds for every $b\in B$, so
\[
\boxed{B\subseteq X+(A-A),\qquad |X|\le K.}
\]
There are at most $K$ such translates, where $|X|$ is an integer even if the constant $K$ is not.

\tools Finite injections, sumsets, and maximal disjoint collections.

\leads A small sumset forces more than a numerical inequality: it lets us cover a second set by a small number of translates of a difference set. That is one way additive structure becomes visible.
\end{solution}

\begin{solution}{pb:returnbound}
\idea Split a walk of $2t$ steps into two independent walks of $t$ steps. Symmetry makes the chance of returning to the identity a sum of squares. The Cauchy--Schwarz inequality then explains why the identity is the most likely location after an even number of steps.

\writeup Let $\mu(g)$ denote the one-step distribution and assume $\mu(g)=\mu(g^{-1})$. Let $\xi_1,\ldots,\xi_t$ be independent increments with law $\mu$, and $X_t=\xi_1\cdots\xi_t$. Because the increments are independent and $\mu$ is symmetric, the reverse-inverse product
\[
X_t^{-1}=\xi_t^{-1}\cdots\xi_1^{-1}
\]
has the same distribution as $X_t$. Hence $p_t(g^{-1})=p_t(g)$ for every $g$.

\emph{(a) A return probability is a squared norm.}
The first $t$ increments and the next $t$ increments are independent, and each block has distribution $p_t$. Conditioning on the first block gives
\[
p_{2t}(g)=\sum_{h\in G}p_t(h)p_t(h^{-1}g).
\]
For $g=e$, symmetry turns this into
\[
\boxed{p_{2t}(e)=\sum_h p_t(h)^2.}
\]
For arbitrary $g$, apply Cauchy--Schwarz:
\[
\begin{aligned}
p_{2t}(g)
&=\sum_h p_t(h)p_t(h^{-1}g)\\
&\le\left(\sum_h p_t(h)^2\right)^{1/2}
 \left(\sum_h p_t(h^{-1}g)^2\right)^{1/2}\\
&=p_{2t}(e).
\end{aligned}
\]
The second sum of squares equals the first because $h\mapsto h^{-1}g$ is a bijection. Finally, apply Cauchy--Schwarz to the constant function $1$ and the probability mass function $p_t$:
\[
1=\left(\sum_hp_t(h)\right)^2
\le |G|\sum_hp_t(h)^2.
\]
Thus $\boxed{p_{2t}(e)\ge1/|G|}$.

\emph{(b) Even-time returns cannot increase.}
For a real function $f$ on $G$, write $\|f\|_2=(\sum_g f(g)^2)^{1/2}$. The one-step update is an average of translates:
\[
p_{t+1}(g)=\sum_{s\in G}\mu(s)p_t(gs^{-1}).
\]
Every translate has the same $\ell^2$ norm as $p_t$, so the triangle inequality gives
\[
\|p_{t+1}\|_2\le\sum_s\mu(s)\|p_t\|_2=\|p_t\|_2.
\]
Use (a) at times $t$ and $t+1$:
\[
\boxed{p_{2t+2}(e)=\|p_{t+1}\|_2^2
\le\|p_t\|_2^2=p_{2t}(e).}
\]

\emph{(c) Odd times behave differently.}
On $\mathbb Z/2\mathbb Z$, move by $1$ at each step. This step distribution is symmetric because $1=-1$ in this group, but at every odd time the walk is at $1$, so $p_{2t+1}(0)=0$.

For any symmetric walk and any $g$, split $2t+1$ increments into blocks of lengths $t$ and $t+1$. The same convolution and Cauchy--Schwarz arguments give
\[
p_{2t+1}(g)\le\|p_t\|_2\|p_{t+1}\|_2
\le\|p_t\|_2^2
=\boxed{p_{2t}(e)}.
\]
Notice that the inequality compares an odd-time probability with the preceding even-time return probability; it does not assert monotonicity of odd-time returns.

\emph{(d) A cycle walk.}
Put $\omega=e^{2\pi i/n}$. The finite geometric sum from Volume~1 gives, for integers $x,j$,
\[
\frac1n\sum_{k=0}^{n-1}\omega^{k(x-j)}=
\begin{cases}1,&x\equiv j\pmod n,\\0,&x\not\equiv j\pmod n.\end{cases}
\]
Let $S_s$ be the ordinary integer sum of $s$ independent equally likely signs. Taking the expectation of this finite identity at $x=S_s$ gives
\[
p_s(j)=\frac1n\sum_{k=0}^{n-1}\omega^{-kj}
\left(\frac{\omega^k+\omega^{-k}}2\right)^s
=\frac1n\sum_{k=0}^{n-1}\omega^{-kj}\lambda_k^s,
\quad \lambda_k=\cos(2\pi k/n).
\]
Independence factors the expectation of the product into $s$ one-step averages. At $s=0$ the empty product is $1$, so every zeroth power in this formula is interpreted as $1$. Pairing conjugate summands replaces $\omega^{-kj}$ by $\cos(2\pi kj/n)$ in the real sum. In particular,
\[
p_{2t}(0)=\frac1n\sum_{k=0}^{n-1}\lambda_k^{2t}.
\]
All terms are non-negative, and the $k=0$ term equals $1/n$, giving the lower bound of (a) directly. As $|\lambda_k|\le1$, each term $\lambda_k^{2t}$ is non-increasing in $t$, which proves (b) again. Also
\[
p_{2t}(j)=\frac1n\sum_k\lambda_k^{2t}\cos(2\pi kj/n)
\le\frac1n\sum_k\lambda_k^{2t}=p_{2t}(0),
\]
which checks the remaining claim of (a). The formula works even when $n$ is even and the walk is periodic.

\tools Finite independence and Cauchy--Schwarz from Volume~1; its roots-of-unity geometric sum. The convolution and cycle formulas were derived here.

\leads The identity $p_{2t}(e)=\|p_t\|_2^2$ is a first connection between a random walk and a norm. It is the bridge from finite-group symmetry to quantitative mixing estimates in Volume~4.
\end{solution}

\section*{Sources and further reading}

\begin{itemize}
\item M.~Artin, \emph{Algebra}, 2nd ed., Pearson, 2011, Chapters~2 and~6. Groups and geometric symmetry, including dihedral groups and finite rotation groups.
\item M.~A.~Armstrong, \emph{Groups and Symmetry}, Springer, 1988. Geometric examples of permutations, matrix groups and rotations of the regular solids.
\item N.~Carter, \emph{Visual Group Theory}, Mathematical Association of America, 2009. Cayley diagrams and multiplication tables as pictures, in the spirit of Figure~\ref{fig:symmetries}.
\item D.~S.~Dummit and R.~M.~Foote, \emph{Abstract Algebra}, 3rd ed., Wiley, 2004, Chapters~1--3. Groups, subgroups and quotient groups, with additional examples and exercises.
\item W.~Burnside, \emph{Theory of Groups of Finite Order}, 2nd ed., Cambridge University Press, 1911 (Dover reprint, 1955). An early systematic treatment of finite groups.
\item T.~Rokicki, H.~Kociemba, M.~Davidson and J.~Dethridge, \emph{The diameter of the Rubik's cube group is twenty}, SIAM J.\ Discrete Math.\ 27 (2013), 1082--1105 (the computation was announced in 2010).
\end{itemize}
